% Supporting Information: detailed methods for
% "Income inequality increases optimal carbon price" (Papargyri & Caldeira).
%
% Written from methods/METHODS_DETAIL.md in the model repository. That document is the
% implementation-level reference (file and function names, floating-point
% evaluation order, regression tests); this one states the same methods as
% mathematics, for the Supporting Information of the paper.
%
% Lines beginning "% AUTHOR NOTE" mark points the authors should decide on or
% verify before submission. They do not appear in the compiled PDF.

\documentclass[11pt]{article}

\usepackage[margin=1in]{geometry}
\usepackage{amsmath}
\usepackage{amssymb}
\usepackage{booktabs}
\usepackage{array}
\usepackage{longtable}
\usepackage[hidelinks]{hyperref}

\renewcommand{\thesection}{S\arabic{section}}
\renewcommand{\theequation}{S\arabic{equation}}
\renewcommand{\thetable}{S\arabic{table}}

% Wider number boxes in the table of contents for "S10"-style numbers
% (article.cls definitions with larger widths).
\makeatletter
\renewcommand*\l@section[2]{%
  \ifnum \c@tocdepth >\z@
    \addpenalty\@secpenalty
    \addvspace{1.0em \@plus\p@}%
    \setlength\@tempdima{2.4em}%
    \begingroup
      \parindent \z@ \rightskip \@pnumwidth
      \parfillskip -\@pnumwidth
      \leavevmode \bfseries
      \advance\leftskip\@tempdima
      \hskip -\leftskip
      #1\nobreak\hfil
      \nobreak\hb@xt@\@pnumwidth{\hss #2\kern-\p@\kern\p@}\par
    \endgroup
  \fi}
\renewcommand*\l@subsection{\@dottedtocline{2}{2.4em}{3.4em}}
\makeatother
\setlength{\emergencystretch}{2em}

\newcommand{\clip}{\operatorname{clip}}
\newcommand{\dd}{\mathrm{d}}
\newcommand{\Lor}{\mathcal{L}}
\newcommand{\etat}{\eta_{\tau}}
\newcommand{\yg}{y_{\mathrm{g}}}
\newcommand{\yd}{y_{\mathrm{d}}}
\newcommand{\yn}{y_{\mathrm{n}}}
\newcommand{\ya}{y_{\mathrm{a}}}
\newcommand{\Ecum}{E_{\mathrm{cum}}}
\newcommand{\Eland}{E_{\mathrm{land}}}
\newcommand{\epot}{e_{\mathrm{pot}}}
\newcommand{\Omb}{\Omega_{\mathrm{b}}}
\newcommand{\yref}{y_{\mathrm{ref}}}
\newcommand{\tb}{t_{\mathrm{b}}}
\newcommand{\Nq}{N_{q}}
\newcommand{\tc}{t^{\mathrm{c}}}

\title{Supporting Information\\[4pt]
\large Detailed methods for ``Income inequality increases optimal carbon price''}
\author{Lamprini Papargyri and Ken Caldeira}
\date{}

\begin{document}

\maketitle

\noindent
This Supporting Information specifies the model, the numerical procedures, the
design of the simulation experiments, the estimation of the model's exogenous
inputs and the construction of the figures. Together with the configuration
file of each run and the results archived with the code, it is intended to
determine every reported number. A companion document in the code repository
(\texttt{methods/METHODS\_DETAIL.md}) gives the same specification at the level of the
implementation, including the floating-point evaluation order needed to
reproduce results bit for bit.

\tableofcontents

%======================================================================
\section{Overview and reproducibility}
\label{sec:overview}
%======================================================================

The model is a stylized, globally aggregated climate--economy model in the
tradition of DICE~\cite{BarrageNordhaus2024} and COIN~\cite{Caldeira2023},
extended with a continuous income distribution. Output is produced by a
Cobb--Douglas technology and is reduced by climate damage and by the cost of
emission abatement. Damage is distributed across the income distribution
according to a vulnerability exponent $\beta$, and the cost of abatement is
raised by a tax whose progressivity is set by a parameter $\tau$. Warming is
computed from cumulative emissions with a three-mode impulse-response model.
Two control trajectories, the carbon price (equivalently the marginal abatement
cost) and the savings rate, are chosen to maximize the discounted sum of
population-weighted mean utility.

The calculation has three stages, which are reproducible to different degrees.
\begin{enumerate}
  \item \emph{Forward integration} of the model for given control trajectories
  (Sections~\ref{sec:inputs}--\ref{sec:derived}) is deterministic. Replaying
  the published optimal controls through the model reproduces every published
  time series to a relative difference of $3.4\times10^{-13}$ or less.
  \item \emph{The objective} (Section~\ref{sec:objective}) inherits the
  precision of the forward integration.
  \item \emph{The optimization} (Sections~\ref{sec:controls}
  and~\ref{sec:optimization}) uses a deterministic derivative-free algorithm,
  but the path it takes depends on every bit of every objective value, so
  implementations that differ only by round-off end at slightly different
  points. Re-optimizing manuscript runs with code whose objective values
  differ from the original only by round-off reproduced the objective to
  $10^{-13}$ (relative), the control $f$ to about $3\times10^{-6}$, and the
  carbon price in 2030 and 2100 and the warming in 2100 to $2\times10^{-6}$
  (relative).
\end{enumerate}
An independent reimplementation of the model and the optimizer, written only
from the repository specification without access to the model's source code,
reproduced the regression baseline (Section~\ref{sec:verification}) bit for
bit.
% AUTHOR NOTE: the independent reimplementation was written by an AI coding
% agent (2026-10-26). Decide whether and how to disclose this.

All arithmetic is IEEE-754 double precision. The reference results were
produced with Python~3.10.12, NumPy~2.2.6~\cite{Harris2020},
SciPy~1.15.3~\cite{Virtanen2020} and the NLopt optimization
library~\cite{NLopt}.
% AUTHOR NOTE: the NLopt version used for the 2026 production runs was not
% recorded; 2.11.0 is the version of the test environment.

%======================================================================
\section{Conventions and notation}
\label{sec:notation}
%======================================================================

\paragraph{Time.} The model time $t$ is the calendar year. Exogenous time
functions and control times are expressed in relative time $t_r = t - \tb$,
with $\tb = 2020$. The integration runs from $t_0 = 2020$ to $t_{\mathrm{end}}
= 2420$ with time step $\Delta t = 1$~yr.

\paragraph{Units.} Monetary quantities are in 2019 US dollars (the DICE-2023
convention), and income is in dollars per person per year. Emissions are in
tCO$_2$ (cumulative) or tCO$_2$~yr$^{-1}$ (rates), and temperature is in
$^\circ$C above preindustrial.

\paragraph{Per-capita and aggregate quantities.} Lower-case symbols denote
per-capita quantities and upper-case symbols denote totals; for example,
gross output is $Y_{\mathrm{g}} = \yg L$ for population $L$.

\paragraph{Population rank.} $F\in[0,1]$ is the rank in the income
distribution, with $F=0$ the poorest. The Lorenz curve $\Lor(F)$ is the share
of total income received by the poorest fraction $F$ of the population, so
that income at rank $F$ is proportional to $\Lor'(F)$, which has mean one.

\paragraph{Clipping.} We write $\clip(x; a, b) = \min\{\max\{x, a\}, b\}$.
Several quantities are bounded away from zero or one by the small constant
$\varepsilon = 10^{-12}$; these safeguards are collected in
Section~\ref{sec:caveats}.

\paragraph{Principal symbols.} Table~\ref{tab:symbols} lists the symbols used
in the model equations.

\begin{longtable}{>{$}l<{$} p{0.72\textwidth}}
\caption{Principal symbols.}\label{tab:symbols}\\
\toprule
\text{Symbol} & Meaning \\
\midrule
\endfirsthead
\toprule
\text{Symbol} & Meaning \\
\midrule
\endhead
\bottomrule
\endfoot
K & capital stock (\$) \\
A, L & total factor productivity; population (persons) \\
\yg, \yd, \yn & per-capita gross output, output after climate damage, and
  output after damage and abatement cost (\$ person$^{-1}$ yr$^{-1}$) \\
\ya & per-capita abatement expenditure \\
\Omb & aggregate damage fraction implied by the current warming \\
\Omega, \omega_i & aggregate and per-node damage fractions \\
q_i, Q & per-node and aggregate damage ratios carried to the next step \\
f, P & control $f = \log_{10} P$; carbon price $P$ (\$ tCO$_2^{-1}$), equal
  to the marginal abatement cost \\
s & savings rate (control) \\
\mu & abatement fraction \\
\sigma, \xi, \Eland & CO$_2$ intensity of output; ratio of CO$_2$-equivalent
  to industrial CO$_2$ emissions; land-use emissions \\
\theta_1, \theta_2 & backstop price (\$ tCO$_2^{-1}$); abatement cost exponent \\
\Ecum, H_k & cumulative emissions; temperature contributions of the three
  impulse-response modes \\
\Delta T & warming above preindustrial \\
G & Gini coefficient of pre-damage income \\
\Lor, \Lor_b & Lorenz curve; base (empirical) Lorenz curve \\
F_i, w_i & quadrature nodes and weights in population rank, $i = 1,\dots,\Nq$ \\
\eta, \rho & elasticity of marginal utility; pure rate of time preference \\
\tau, \etat & tax progressivity parameter; exponent of the tax schedule \\
\beta & damage vulnerability exponent \\
\Delta U & utility loss parameter of the tax schedule \\
U & population-weighted mean utility, $U = L\,\bar u$ \\
\varepsilon & numerical safeguard constant, $10^{-12}$ \\
\end{longtable}

%======================================================================
\section{Exogenous inputs}
\label{sec:inputs}
%======================================================================

\subsection{Scalar parameters}

Each run is specified by a configuration that fixes all parameters; there are
no implicit defaults. Table~\ref{tab:scalars} lists the scalar parameters.
Only $\eta$, $\rho$, $\tau$ and $\beta$ vary between the manuscript runs
(Section~\ref{sec:design}).

\begin{table}[ht]
\centering
\caption{Scalar parameters of the manuscript runs.}
\label{tab:scalars}
\small
\begin{tabular}{l l l}
\toprule
Symbol & Value & Meaning \\
\midrule
$\alpha$ & 0.3 & capital share of output \\
$\delta$ & 0.1 yr$^{-1}$ & capital depreciation rate \\
$\psi_1$ & 0 $^\circ$C$^{-1}$ & linear damage coefficient \\
$\psi_2$ & 0.003467 $^\circ$C$^{-2}$ & quadratic damage coefficient \\
$\Omega_{\max}$ & 0.8 & upper bound on the damage fraction \\
$\eta$ & 1.5 (varied) & elasticity of marginal utility \\
$\rho$ & 0.025 yr$^{-1}$ (varied) & pure rate of time preference \\
$\theta_2$ & 2.6 & abatement cost exponent \\
$\mu_{\max}$ & $10^{20}$ & upper bound on $\mu$ (effectively none) \\
$K_0$ & $2.95\times10^{14}$ \$ & capital in 2020 \\
$\yref$ & 17\,210.79 \$ person$^{-1}$ yr$^{-1}$ & reference income of the
  damage distribution \\
$\beta$ & 0, 0.5, 1 (varied) & damage vulnerability exponent \\
$\tau$ & 0, 0.5, 0.998 (varied) & tax progressivity, $0\le\tau<1$ \\
$\Nq$ & 32 & number of quadrature nodes \\
\bottomrule
\end{tabular}
\end{table}

$\mu > 1$ is allowed and represents net-negative industrial emissions. The
model also supports emission and consumption pulses at a given year
(Section~\ref{sec:pulses}); their amplitudes are zero in all manuscript runs.

\subsection{Time functions}

Exogenous time-dependent inputs are given by one of the following functional
forms, evaluated at $t_r = t-\tb$:
\begin{align}
  &\text{exponential:} && g(t_r) = a\,e^{b t_r} + c, \label{eq:tf-exp}\\
  &\text{double exponential:} && g(t_r) = g_0\left[\phi\, e^{b_1 t_r}
     + (1-\phi)\,e^{b_2 t_r}\right], \label{eq:tf-dexp}\\
  &\text{Gompertz:} && g(t_r) = g_\infty \exp\!\left[\ln\!\left(\frac{g_0}{g_\infty}\right)
     e^{c\, t_r}\right], \label{eq:tf-gompertz}\\
  &\text{logistic:} && g(t_r) = \frac{g_\infty}{1 + \left(g_\infty/g_0 - 1\right)e^{-b t_r}}, \\
  &\text{piecewise linear:} && \text{linear interpolation between given points,
     constant beyond the end points,}
\end{align}
or a constant. The Gompertz form has $g(0) = g_0$ and $g\to g_\infty$ for
$c<0$. Table~\ref{tab:timefunctions} gives the time functions of the
manuscript runs; Section~\ref{sec:estimation} gives their provenance.

\begin{table}[ht]
\centering
\caption{Exogenous time functions of the manuscript runs ($t_r = t - 2020$).}
\label{tab:timefunctions}
\small
\begin{tabular}{l l p{0.5\textwidth}}
\toprule
Input & Form & Parameters \\
\midrule
$A(t)$ & Gompertz & $g_0 = 739.619$, $g_\infty = 4\,906\,815.811$,
  $c = -0.0015$ yr$^{-1}$ \\
$L(t)$ (persons) & Gompertz & $g_0 = 7.752\times10^9$,
  $g_\infty = 1.0825\times10^{10}$, $c = -0.03133$ yr$^{-1}$ \\
$\sigma(t)$ (tCO$_2$ \$$^{-1}$) & Gompertz & $g_0 = 2.91355\times10^{-4}$,
  $g_\infty = 4.4681\times10^{-5}$, $c = -0.008164$ yr$^{-1}$ \\
$\theta_1(t)$ (\$ tCO$_2^{-1}$) & double exponential & $g_0 = 6035.6$,
  $\phi = 0.906$, $b_1 = -0.224$, $b_2 = -0.00112$ yr$^{-1}$ \\
$G(t)$ & exponential & $a = 0.19517$, $b = -0.01058$ yr$^{-1}$,
  $c = 0.37796$ \\
$\xi(t)$ & exponential & $a = 1.4$, $b = -0.002085$ yr$^{-1}$, $c = 0$ \\
$\Eland(t)$ (tCO$_2$ yr$^{-1}$) & exponential & $a = 5.9\times10^9$,
  $b = -0.02107$ yr$^{-1}$, $c = 0$ \\
\bottomrule
\end{tabular}
\end{table}

The Gini trajectory gives $G(2020) = 0.57313$ and $G(\infty) = 0.37796$. In
the runs that compare the model with DICE-2023 (Section~\ref{sec:design}),
$G(t)\equiv 0$.

\subsection{Climate response and initial state}
\label{sec:climate-constants}

Warming responds to cumulative CO$_2$ emissions through the impulse response
of the median model of Ricke and Caldeira~\cite{RickeCaldeira2014}. The warming
at time $t$ caused by a unit emission at time $t'$ is
\begin{equation}
  R(t - t') = R_\infty + \sum_{k=1}^{3} a_k\, e^{-(t-t')/\zeta_k},
  \qquad R_\infty = -(a_1+a_2+a_3),
  \label{eq:impulse}
\end{equation}
with $a = (-0.629,\ 0.203,\ -0.0521)\times10^{-12}$ $^\circ$C~tCO$_2^{-1}$
and $\zeta = (2.241,\ 35.750,\ 97.180)$~yr, so that $R_\infty =
4.781\times10^{-13}$ $^\circ$C~tCO$_2^{-1}$ and $R(0) = 0$.

The 2020 climate state comes from a fit of historical warming,
\begin{equation}
  \Delta T_{\mathrm{hist}}(t) = 1.1126 - 0.07022\,e^{(2020-t)/\zeta_1}
  + 0.21036\,e^{(2020-t)/\zeta_2} - 0.08308\,e^{(2020-t)/\zeta_3}
  \quad (^\circ\mathrm{C}),
\end{equation}
obtained by convolving an exponential emission history (633.5~GtC cumulative
to 2020, as in DICE-2023) with Eq.~\eqref{eq:impulse}. Its constant term is
$R_\infty\Ecum(2020)$, which gives $\Ecum(2020) = 1.1126/R_\infty =
2.327128\times10^{12}$~tCO$_2$. The other three terms are the 2020 mode
temperatures, $H_1 = -0.07022$, $H_2 = 0.21036$ and $H_3 = -0.08308$~$^\circ$C.
The initial warming is $\Delta T(2020) = 1.16966$~$^\circ$C.

%======================================================================
\section{Income distribution}
\label{sec:distribution}
%======================================================================

\subsection{Quadrature over population rank}

Population averages are computed with $\Nq$-point Gauss--Legendre quadrature
mapped to $[0,1]$. If $x_i$ and $W_i$ are the Gauss--Legendre nodes and
weights on $[-1,1]$, the rank nodes and population weights are
\begin{equation}
  F_i = \tfrac{1}{2}(x_i + 1), \qquad w_i = \tfrac{1}{2} W_i,
  \qquad \sum_{i=1}^{\Nq} w_i = 1,
\end{equation}
and the population mean of a quantity $z$ is $\langle z\rangle =
\sum_i w_i z_i$. Node $i$ represents the population interval
$[F_{i-}, F_{i+}]$ with $F_{1-}=0$ and $F_{i+} - F_{i-} = w_i$; by the
Chebyshev--Markov--Stieltjes separation theorem each node lies strictly inside
its interval, so node values also serve as piecewise-constant values on the
intervals. For $\Nq = 32$, $F_1 = 0.001368$, $F_{32} = 0.998632$ and
$w_1 = 0.003509$.

\subsection{Empirical Lorenz curve}
\label{sec:lorenz}

All manuscript runs use a Lorenz curve fitted to the global income
distribution (Section~\ref{sec:lorenz-fit}). The base curve and its derivative
are
\begin{equation}
  \Lor_b(F) = \sum_{k=0}^{3} \gamma_k F^{p_k}, \qquad
  \Lor_b'(F) = \sum_{k=0}^{3} \gamma_k\, p_k\, F^{p_k - 1},
  \label{eq:lorenz-base}
\end{equation}
with exponents and weights
\begin{gather*}
  p = (1.500036,\ 4.367440,\ 14.072005,\ 135.059674),\\
  \gamma_1 = 0.3776187268483524,\quad \gamma_2 = 0.3671247620949191,\\
  \gamma_3 = 0.09538538350961864,\quad \gamma_0 = 1-\gamma_1-\gamma_2-\gamma_3 .
\end{gather*}
Its Gini coefficient is
\begin{equation}
  G_b = 1 - 2\sum_{k=0}^{3}\frac{\gamma_k}{p_k+1} = 0.681279513926459.
\end{equation}
A Lorenz curve with any Gini coefficient $0\le G\le G_b$ is obtained by mixing
the base curve linearly with perfect equality:
\begin{equation}
  \Lor(F) = (1-a)\,F + a\,\Lor_b(F), \qquad
  \Lor'(F) = (1-a) + a\,\Lor_b'(F), \qquad a = G/G_b,
  \label{eq:lorenz-mix}
\end{equation}
which has Gini coefficient exactly $G$. A Gini coefficient above $G_b$ is
treated as a configuration error. For $G = 0$, $\Lor'(F) \equiv 1$.

The density $\Lor_b'$ is steep at the top of the distribution ($p_3 = 135$),
so the quadrature does not integrate it exactly: the computed mean
$\sum_i w_i \Lor'(F_i) - 1$ is $3.2\times10^{-7}$ at $G = 0.3$,
$6.05\times10^{-7}$ at $G = 0.573$ and $7.2\times10^{-7}$ at $G = G_b$. The
model does not renormalize; this error is part of the published results.

The model also supports a Pareto Lorenz curve, $\Lor'(F) = (1 - 1/a)(1-F)^{-1/a}$
with $a = (1 + 1/G)/2$. It is not used for any result: its density is singular
at $F\to1$, and with 32 nodes the quadrature loses 2--10\% of income.

%======================================================================
\section{Climate and economy in one time step}
\label{sec:step}
%======================================================================

At each time $t$ the model computes the tendencies of the state variables
$(K, \Ecum, H_1, H_2, H_3)$ from the state, the exogenous inputs, the controls
$f(t)$ and $s(t)$, and two damage-distribution ratios carried from the
previous step: the per-node ratios $q_i$ and the aggregate ratio $Q$. The
calculation proceeds in the order of the following subsections.

\subsection{Temperature}

\begin{equation}
  \Delta T = R_\infty \Ecum + H_1 + H_2 + H_3 .
\end{equation}

\subsection{Climate damage with a lagged distribution}
\label{sec:damage}

The aggregate damage fraction implied by the current warming is
\begin{equation}
  \Omb = \min\left\{\Omega_{\max},\ \psi_1\Delta T + \psi_2 \Delta T^2\right\}.
  \label{eq:damage-base}
\end{equation}
The damage applied in this step uses the current $\Omb$ but the distribution
of damage across incomes computed at the end of the previous step
(Section~\ref{sec:damage-update}):
\begin{equation}
  \tilde\Omega = \clip(Q\,\Omb;\ 0,\ 1-\varepsilon), \qquad
  \tilde\omega_i = \clip(q_i\,\Omb;\ 0,\ 1-\varepsilon).
  \label{eq:damage-lagged}
\end{equation}
In the first step $q_i = Q = 1$, i.e.\ damage is uniform across incomes.

\subsection{Production}

\begin{equation}
  \yg = A\left(\frac{K}{L}\right)^{\alpha}, \qquad
  \yd = (1-\tilde\Omega)\,\yg ,
\end{equation}
with $\yg = 0$ if $K\le0$. Climate damage per capita is $\tilde\Omega\,\yg$.
If $\yg\le\varepsilon$ or $\yd\le\varepsilon$ the economy has collapsed and
the step is handled as described in Section~\ref{sec:collapse}.

\subsection{Abatement}
\label{sec:abatement}

The control sets the carbon price, $P = 10^{f}$, which equals the marginal
abatement cost. Potential (unabated) emissions per capita, including land-use
emissions, are
\begin{equation}
  \epot = \sigma\,\yg\,\xi + \Eland/L .
\end{equation}
Abatement follows the DICE cost curve, whose marginal cost
$\theta_1\mu^{\theta_2-1}$ is set equal to $P$:
\begin{equation}
  \mu = \min\left\{\mu_{\max},\ \left(\frac{P}{\theta_1}\right)^{1/(\theta_2-1)}\right\},
  \qquad
  \ya = \min\left\{\frac{\theta_1}{\theta_2}\,\epot\,\mu^{\theta_2},\
  (1-\varepsilon)\,\yd\right\},
  \label{eq:abatement}
\end{equation}
if $\epot > \varepsilon$ and $P > \varepsilon$; otherwise $\mu = \ya = 0$.
Abatement expenditure is paid from damaged output:
\begin{equation}
  \lambda = \ya/\yd, \qquad \yn = (1-\lambda)\,\yd .
\end{equation}
Land-use emissions are abated at the same fraction $\mu$ as industrial
emissions.

\subsection{Emission and consumption pulses}
\label{sec:pulses}

For impulse experiments, an emission pulse $X_E$ (tCO$_2$) or a consumption
pulse $X_C$ (\$) can be added in the step $[t, t+\Delta t)$ that contains a
given year. They contribute $E_{\mathrm{extra}} = X_E/\Delta t$ to the
emission rate and $c_{\mathrm{extra}} = X_C/(L\Delta t)$ to per-capita income
in that step, and zero otherwise. Both are zero in all manuscript runs.

\subsection{Pre-tax income distribution}

Pre-tax income at node $i$ is gross income at that rank, reduced by the
node's lagged damage fraction:
\begin{equation}
  y_{\mathrm{pre},i} = \yg\,\Lor'(F_i)\,(1-\tilde\omega_i),
  \qquad \hat y_i = y_{\mathrm{pre},i}/\yd ,
  \label{eq:pretax}
\end{equation}
where $\Lor'$ uses the current $G(t)$, and $\hat y_i$ is pre-tax income
normalized by mean damaged output. The mean $\langle \hat y\rangle$ equals one
only approximately, because of the quadrature error of
Section~\ref{sec:lorenz} and because the lagged node damages $\tilde\omega_i$
were normalized with the previous step's post-tax incomes
(Section~\ref{sec:damage-update}), whereas $\tilde\Omega$ is applied in
aggregate. The model uses these quantities as they are.

\subsection{Progressive taxation}
\label{sec:tax}

The abatement expenditure is financed by a tax on post-damage income. The per
capita revenue required is $T = \ya - c_{\mathrm{extra}}$ (so a consumption
pulse reduces the tax), or $\hat T = T/\yd$ in normalized units.

\paragraph{Tax schedule.} The tax belongs to the family of equal-utility-loss
schedules: every person loses the same amount $\Delta U \ge 0$ of utility
measured with an isoelastic function of elasticity $\etat$,
$u_\tau(c) = c^{1-\etat}/(1-\etat)$ (or $\ln c$ for $\etat = 1$). The
progressivity parameter $\tau\in[0,1)$ sets
\begin{equation}
  \etat = 1 + \frac{\tau}{1-\tau}\,(\eta - 1).
  \label{eq:eta-tau}
\end{equation}
$\tau = 0$ gives $\etat = 1$, a flat tax rate; $\tau = 0.5$ gives
$\etat = \eta$, equal utility loss under the welfare function; and
$\tau\to1$ gives $\etat\to\infty$, a tax that falls on the highest incomes.
For $\eta = 1.5$ and $\tau = 0.998$, $\etat = 250.5$.

Post-tax normalized income $c(\hat y;\Delta U)$ solves
$u_\tau(\hat y) - u_\tau(c) = \Delta U$:
\begin{equation}
  c(\hat y;\Delta U) =
  \begin{cases}
    \hat y\,e^{-\Delta U}, & \etat = 1,\\[4pt]
    \left[\hat y^{\,1-\etat} + (\etat-1)\,\Delta U\right]^{-1/(\etat-1)}, & \etat > 1,\\[4pt]
    \left[\max\left\{\hat y^{\,1-\etat} - (1-\etat)\,\Delta U,\ 0\right\}\right]^{1/(1-\etat)}, & \etat < 1.
  \end{cases}
  \label{eq:schedule}
\end{equation}
The case $\etat = 1$ is applied when $|\etat - 1| < \varepsilon$. For
$\etat>1$ the schedule is evaluated in logarithms,
\begin{equation}
  \ln c = \ln\hat y - \frac{1}{\etat-1}\,
  \ln\!\left[1 + \exp\!\left(\ln(\etat-1) + \ln\Delta U + (\etat-1)\ln\hat y\right)\right],
\end{equation}
with $\ln\Delta U$ as the unknown. For strongly progressive taxes $\Delta U$
is extremely small (for $\tau = 0.998$, $\ln\Delta U$ lies between about
$-670$ and $-460$), and working with $\ln\Delta U$ avoids an underflow that
would otherwise switch the tax off.

\paragraph{Revenue condition.} If $T \le \varepsilon$ there is no tax. For a
flat tax ($|\etat - 1|<\varepsilon$), every income is reduced by the fraction
$r = \min\{\hat T, 1-\varepsilon\}$, i.e.\ $\Delta U = -\ln(1-r)$. This closed
form assumes $\langle\hat y\rangle = 1$, so the revenue raised differs from
$\hat T$ by the small relative error discussed above. Otherwise
$x = \ln\Delta U$ is the root of the revenue residual
\begin{equation}
  g(x) = \sum_{i} w_i\left[\hat y_i - c\left(\hat y_i; e^{x}\right)\right] - \hat T ,
  \label{eq:revenue}
\end{equation}
found with Brent's method~\cite{Brent1973} on the bracket
\begin{equation}
  x_{\mathrm{lo}} = -\left|(\etat-1)\ln\max_i \hat y_i\right| - 50,
  \qquad x_{\mathrm{hi}} = \ln 10^{10},
\end{equation}
with an absolute tolerance of $10^{-9}$ in $x$, a relative tolerance of
$4\cdot2^{-52}$ and at most 100 iterations. If $g(x_{\mathrm{lo}})\ge0$, the
target is negligible and no tax is levied. If $g(x_{\mathrm{hi}})<0$, the
revenue cannot be raised and the control trajectory is infeasible
(Section~\ref{sec:nlopt}). A solution whose post-tax incomes are not finite,
or decrease with rank by more than $10^{-9}$, is also treated as a failure.

Post-tax income is $y_{\mathrm{n},i} = c(\hat y_i;\Delta U)\,\yd$.

\paragraph{Tax-rate diagnostics.} The average and marginal tax rates at node
$i$ are reported but do not affect the dynamics:
\begin{equation}
  t^{\mathrm{avg}}_i = 1 - \frac{y_{\mathrm{n},i}}{y_{\mathrm{pre},i}}, \qquad
  t^{\mathrm{marg}}_i = 1 - \frac{\partial c}{\partial\hat y}
  = 1 - \left(\frac{y_{\mathrm{n},i}}{y_{\mathrm{pre},i}}\right)^{\etat},
\end{equation}
with the ratio $y_{\mathrm{n},i}/y_{\mathrm{pre},i}$ set to one where
$y_{\mathrm{pre},i}\le\varepsilon$, and the marginal rate set to one where that
ratio is at most $\varepsilon$.

\subsection{Utility}

Every person saves the same fraction $s$ of post-tax income. Consumption and
utility at node $i$ are
\begin{equation}
  c_i = (1-s)\,y_{\mathrm{n},i}, \qquad
  u_i = \begin{cases}
    \dfrac{\max\{c_i,\varepsilon\}^{1-\eta}}{1-\eta}, & \eta\ne1,\\[8pt]
    \ln\max\{c_i,\varepsilon\}, & \eta = 1,
  \end{cases}
\end{equation}
and the population-weighted mean utility is
\begin{equation}
  U = L\,\bar u, \qquad \bar u = \sum_i w_i u_i .
\end{equation}

\subsection{Damage distribution for the next step}
\label{sec:damage-update}

Damage fractions are distributed across incomes by a power law in post-tax
income,
\begin{equation}
  \omega_i^{*} = \clip\!\left(\Omb\left(\frac{y_{\mathrm{n},i}}{\yref}\right)^{-\beta};\
  0,\ 1-\varepsilon\right),
  \label{eq:damage-powerlaw}
\end{equation}
with $\omega_i^{*} = \Omb$ for $\beta = 0$ (applied when $|\beta|<\varepsilon$)
and $\omega_i^{*} = 1-\varepsilon$ where $y_{\mathrm{n},i}\le\varepsilon$. The
fractions are then rescaled so that the income-weighted mean damage fraction
equals $\Omb$:
\begin{equation}
  \omega_i = \clip\!\left(\omega_i^{*}\,
  \frac{\Omb\sum_j w_j y_{\mathrm{n},j}}{\sum_j w_j\,y_{\mathrm{n},j}\,\omega_j^{*}};\
  0,\ 1-\varepsilon\right),
  \qquad
  \Omega = \frac{\sum_i w_i\,y_{\mathrm{n},i}\,\omega_i}{\sum_i w_i\,y_{\mathrm{n},i}} .
  \label{eq:damage-rescale}
\end{equation}
Thus $\beta$ redistributes damage across incomes without changing its
aggregate, and $\yref$ cancels except through the clipping bounds. For
$\beta = 1$, absolute damage $\omega_i\,y_{\mathrm{n},i}$ is the same for every
person; for $\beta > 0$, damage falls disproportionately on lower incomes.
(The rescaling is skipped in an optional variant, not used for any result, in
which the aggregate damage depends on the income distribution.)

The ratios carried to the next step are
\begin{equation}
  q_i = \omega_i/\Omb, \qquad Q = \Omega/\Omb,
\end{equation}
or zero if $\Omb\le\varepsilon$. The next step multiplies them by its own
$\Omb$ (Eq.~\ref{eq:damage-lagged}): the temperature dependence of damage is
current, while its distribution across incomes lags by one time step.

\subsection{Emissions, capital and climate tendencies}

Net emissions per capita and the tendencies of the state are
\begin{align}
  e &= (1-\mu)\,\epot , \\
  \frac{\dd K}{\dd t} &= s\,\yn L - \delta K , \label{eq:capital}\\
  \frac{\dd\Ecum}{\dd t} &= e\,L + E_{\mathrm{extra}} , \\
  \frac{\dd H_k}{\dd t} &= a_k\,\frac{\dd\Ecum}{\dd t} - \frac{H_k}{\zeta_k},
  \qquad k = 1,2,3 . \label{eq:modes}
\end{align}
Investment uses aggregate net output $\yn$, not the mean of the node incomes.
Eqs.~\eqref{eq:modes} are the state-space form of the impulse response
Eq.~\eqref{eq:impulse}: together with $\Delta T = R_\infty\Ecum + \sum_k H_k$
they reproduce the warming of an arbitrary emission history. Emissions are
negative when $\mu > 1$.

\subsection{Economic collapse}
\label{sec:collapse}

If $\yg\le\varepsilon$ or $\yd\le\varepsilon$, the step is treated as an
economic collapse: abatement, taxes, emissions and all per-node incomes and
damages are set to zero, $\dd K/\dd t = -\delta K$, and the mean utility is
assigned the sentinel value $\bar u = U = -10^{30}$ (not multiplied by $L$).
The climate tendencies and the carried ratios are then computed as above.
This branch is not reached in any of the reference runs.

%======================================================================
\section{Time integration}
\label{sec:integration}
%======================================================================

The model is integrated on the grid $t_n = t_0 + n\,\Delta t$,
$n = 0,\dots,N-1$, which for 2020--2420 with $\Delta t = 1$~yr has $N = 401$
points. The initial state is $K = K_0$ and the 2020 climate state of
Section~\ref{sec:climate-constants}, with $q_i = Q = 1$. At each grid point
the tendencies are evaluated from the state at $t_n$, all diagnostics are
recorded, and the state is advanced with the forward Euler method:
\begin{equation}
  K_{n+1} = K_n + \Delta t\,\dot K_n, \quad
  E_{\mathrm{cum},n+1} = \max\{0,\ E_{\mathrm{cum},n} + \Delta t\,\dot E_{\mathrm{cum},n}\}, \quad
  H_{k,n+1} = H_{k,n} + \Delta t\,\dot H_{k,n}.
\end{equation}
The recorded state at $t_n$ is the state before that step's update; the update
computed at the final grid point is not used. Cumulative emissions are kept
non-negative, so net-negative emissions cannot cool the climate below that of
zero cumulative emissions. No other state is constrained.

%======================================================================
\section{Derived quantities}
\label{sec:derived}
%======================================================================

\paragraph{Aggregates.} Totals are per-capita values times population:
$Y_{\mathrm{g}} = \yg L$, $Y_{\mathrm{d}} = \yd L$, $Y_{\mathrm{n}} = \yn L$,
$E = eL$. Mean consumption and savings per capita are $(1-s)\yn$ and $s\yn$.

\paragraph{Carbon price.} The reported marginal abatement cost is
$\theta_1\mu^{\theta_2-1}$, which equals $P = 10^f$ except where $\mu$ was
capped or set to zero. In the figures, the optimal carbon price in a year is
this value at that grid point (the grid is annual, so no interpolation is
needed).

\paragraph{Consumption discount rate.} The model reports the consumption
discount rate given by the marginal-utility-weighted Ramsey rule. With
$c_{j,n} = \max\{(1-s_n)\,y_{\mathrm{n},j}(t_n),\ \varepsilon\}$ and
$M_n = \sum_j w_j\,c_{j,n}^{-\eta}$,
\begin{equation}
  r_n = e^{\rho\Delta t}\,\frac{M_n}{M_{n+1}} - 1,
  \qquad r_{N-1} = r_{N-2}.
  \label{eq:ramsey}
\end{equation}
This is the discrete form of $r = \rho - \frac{\dd}{\dd t}\ln\sum_j w_j\,
u'(c_j)$, i.e.\ $\rho$ plus $\eta$ times the growth rate of consumption
averaged with weights proportional to marginal utility. For a static
distribution it reduces to the representative-agent rule $\rho + \eta g$.
% AUTHOR NOTE: the r_consumption columns in the published results CSVs were
% computed with the earlier representative-agent rule,
% r = exp(rho) * (cbar_{n+1}/cbar_n)^eta - 1, with cbar the mean consumption;
% Eq. (ramsey) differs from them by 1-2% typically and up to ~50% at single
% points. No figure uses this column (Figs 5 and S3 use r_mu; see the r_mu subsection).

%======================================================================
\section{Objective}
\label{sec:objective}
%======================================================================

The objective is the discounted integral of population-weighted mean utility,
\begin{equation}
  J = \int_{t_0}^{t_{\mathrm{end}}} e^{-\rho(t-t_0)}\,U(t)\,\dd t
    \approx \sum_{n=0}^{N-2}\frac{\Delta t}{2}
    \left[e^{-\rho(t_n-t_0)}U_n + e^{-\rho(t_{n+1}-t_0)}U_{n+1}\right],
  \label{eq:objective}
\end{equation}
evaluated with the trapezoidal rule on the integration grid. For the
manuscript configuration $J\approx -6.5\times10^{9}$.

%======================================================================
\section{Control trajectories}
\label{sec:controls}
%======================================================================

Each control, $f(t)$ and $s(t)$, is represented by its values at $n$ control
times and interpolated between them.

\subsection{Control times}
\label{sec:control-times}

Control times are concentrated early in the century, where the optimal
controls change fastest. In relative time, with $t_0 = 0$ and
$t_{\mathrm{end}} = 400$:
\begin{enumerate}
  \item One wind-down point is used if $n\le10$ and two otherwise
  ($n_{\mathrm{end}}$). The wind-down period begins at
  $t_{\mathrm{join}} = t_{\mathrm{end}} - 4/\delta$ (360, i.e.\ 2380, for
  $\delta = 0.1$~yr$^{-1}$), four e-folding times of capital depreciation
  before the end of the run.
  \item The remaining $n_1 = n - n_{\mathrm{end}}\ge3$ points span
  $[t_0, t_{\mathrm{join}}]$ with $m = n_1 - 1$ geometrically growing
  intervals. The first interval is $\Delta t$, and the growth factor
  $1 + r_0$ solves
  \begin{equation}
    \sum_{i=0}^{m-1}\Delta t\,(1+r_0)^i = t_{\mathrm{join}} - t_0 ,
  \end{equation}
  found with Brent's method on $r_0\in(-0.5, 100)$. If the bracket contains no
  root, $r_0 = (t_{\mathrm{join}}-t_0)/(m\Delta t) - 1$ is used instead.
  \item The points are $\tc_0 = t_0$ and $\tc_{i+1} = \tc_i + \Delta t
  (1+r_0)^i$, accumulated sequentially, with the last point then set exactly to
  $t_{\mathrm{join}}$.
  \item The wind-down points are $t_{\mathrm{end}}$ if $n_{\mathrm{end}}=1$, or
  $t_{\mathrm{end}} - 1/\delta$ and $t_{\mathrm{end}}$ if $n_{\mathrm{end}}=2$.
\end{enumerate}
Table~\ref{tab:control-times} lists the resulting grids. For $n=4$, the
fallback $r_0$ is used; it does not matter, because the second point is
$t_0 + \Delta t$ for any $r_0$ and the third is set to $t_{\mathrm{join}}$.

\begin{table}[ht]
\centering
\caption{Control times (calendar years) of the five refinement iterations.}
\label{tab:control-times}
\small
\begin{tabular}{r l p{0.62\textwidth}}
\toprule
$n$ & $r_0$ & Control times \\
\midrule
4 & 179 (fallback) & 2020, 2021, 2380, 2420 \\
6 & 5.744311 & 2020, 2021, 2027.744, 2073.230, 2380, 2420 \\
9 & 1.443996 & 2020, 2021, 2023.444, 2029.417, 2044.015, 2079.693,
  2166.891, 2380, 2420 \\
15 & 0.555777 & 2020, 2021, 2022.556, 2024.976, 2028.742, 2034.600, 2043.715,
  2057.895, 2079.957, 2114.279, 2167.677, 2250.753, 2380, 2410, 2420 \\
24 & 0.236343 & 2020, 2021, 2022.236, 2023.765, 2025.655, 2027.991, 2030.880,
  2034.451, 2038.867, 2044.326, 2051.075, 2059.419, 2069.736, 2082.490,
  2098.259, 2117.755, 2141.859, 2171.660, 2208.504, 2254.055, 2310.372, 2380,
  2410, 2420 \\
\bottomrule
\end{tabular}
\end{table}

\subsection{Interpolation}
\label{sec:pchip}

Between control times, each control follows the shape-preserving piecewise
cubic Hermite interpolant (PCHIP) of Fritsch and
Butland~\cite{FritschCarlson1980,FritschButland1984}, as implemented in
SciPy~1.15. For control points $(\tc_k, v_k)$ with $h_k = \tc_{k+1}-\tc_k$
and secant slopes $m_k = (v_{k+1}-v_k)/h_k$, the derivative at an interior
point is zero if $m_{k-1}$ and $m_k$ differ in sign or either is zero, and
otherwise the weighted harmonic mean
\begin{equation}
  d_k = \frac{w_1 + w_2}{w_1/m_{k-1} + w_2/m_k}, \qquad
  w_1 = 2h_k + h_{k-1},\quad w_2 = h_k + 2h_{k-1}.
\end{equation}
At the first point,
\begin{equation}
  d_0 = \frac{(2h_0 + h_1)\,m_0 - h_0\,m_1}{h_0 + h_1},
\end{equation}
set to zero if its sign differs from that of $m_0$, and to $3m_0$ if $m_0$
and $m_1$ differ in sign and $|d_0| > 3|m_0|$; the last point is treated
symmetrically, and with only two points both derivatives equal the secant
slope. On $[\tc_k,\tc_{k+1}]$ the interpolant is the cubic Hermite
polynomial with values $v_k, v_{k+1}$ and derivatives $d_k, d_{k+1}$. Beyond
the last control time the control is held at its last value. No bounds are
imposed on the interpolated trajectory; the bounds of
Section~\ref{sec:optimization} apply only to the values at the control times.

%======================================================================
\section{Optimization}
\label{sec:optimization}
%======================================================================

\subsection{Refinement schedule}

The optimization proceeds through $K_{\mathrm{it}} = 5$ iterations with
progressively finer control grids and larger evaluation budgets. For an
initial value $a$ and a final value $b$, iteration $k = 1,\dots,K_{\mathrm{it}}$
uses
\begin{equation}
  \operatorname{round}\!\left(1 + (a-1)\,\left[\frac{b-1}{a-1}\right]^{(k-1)/(K_{\mathrm{it}}-1)}\right),
\end{equation}
with ties rounded to even. With $a = 4$ and $b = 24$ control points, the
iterations use 4, 6, 9, 15 and 24 points for each control; with budgets from
2\,000 to 100\,000 objective evaluations, they allow 2\,000, 5\,317, 14\,140,
37\,602 and 100\,000 evaluations.

\subsection{Iterations}

In the first iteration, every control value is set to its initial guess
($f = 2.5$, i.e.\ \$316~tCO$_2^{-1}$; $s = 0.256$), except at the final control
time, 2420, where both are set to zero. In each later iteration, the previous
optimum is interpolated (Section~\ref{sec:pchip}) onto the new control times.
The optimization variables are the values of $f$ and then of $s$ at their
control times, with bounds $f\in[-1, 4]$ (a carbon price between \$0.10 and
\$10\,000~tCO$_2^{-1}$) and $s\in[0, 1]$; the starting point is clipped to the
bounds.

Each iteration maximizes $J$ with the Subplex algorithm~\cite{Rowan1990} as
implemented in the NLopt library~\cite{NLopt} (algorithm \texttt{LN\_SBPLX}).
The only stopping
criteria are an absolute tolerance of $10^{-8}$ on every variable and the
iteration's evaluation budget; relative and absolute tolerances on the
objective are disabled, and all other settings, including the initial step
size, are NLopt's defaults. The result of the final iteration is the reported
optimum. In the nine reference runs the final iteration always terminated on
the variable tolerance.

\subsection{Objective presented to the optimizer}
\label{sec:nlopt}

The optimizer maximizes $10^{-13}J$. If the model cannot be evaluated for a
trial trajectory---because the required tax revenue cannot be raised, the
Gini coefficient exceeds $G_b$, or the tax root solve fails---$J$ is set to
$-10^{30}$, which the optimizer sees as $-10^{17}$. If $J$ is not finite, the
optimizer receives $-10^{30}$ directly.

%======================================================================
\section{Simulation experiments}
\label{sec:design}
%======================================================================

All 14\,801 manuscript runs share the template of
Section~\ref{sec:inputs} and differ only in $\eta$, $\rho$, $\tau$ and $\beta$,
with the exceptions noted. Unless stated otherwise, $\eta = 1.5$ and
$\rho = 0.025$~yr$^{-1}$, and the three tax cases are
$\tau\in\{0, 0.5, 0.998\}$. Table~\ref{tab:design} summarizes the families of
runs. The configuration file of every run is archived with the code, so every
simulation behind the manuscript can be re-run.

\begin{longtable}{p{0.19\textwidth} r p{0.55\textwidth}}
\caption{Families of simulation runs.}\label{tab:design}\\
\toprule
Family & Runs & Design \\
\midrule
\endfirsthead
\toprule
Family & Runs & Design \\
\midrule
\endhead
\bottomrule
\endfoot
Reference cases & 9 &
  $\tau\in\{0, 0.5, 0.998\}\times\beta\in\{0, 0.5, 1\}$. \\
Contour grid & 1\,681 &
  $\tau\in\{k/40: k = 0,\dots,39\}\cup\{0.998\}$ $\times$
  $\beta\in\{k/40: k = 0,\dots,40\}$. \\
Monte Carlo & $2\times1\,500$ &
  For each $\beta\in\{0, 0.5\}$, a separate Latin hypercube
  sample~\cite{McKay1979} of 500 $(\eta,\rho)$ pairs with
  $\eta\in[1,2]$ and $\rho\in[0, 0.05]$~yr$^{-1}$ (rounded to six decimals),
  each run for the three values of $\tau$. \\
Single-parameter $\eta$ & $2\times1\,500$ &
  500 values of $\eta\in[1,2]$, one in each stratum of width 0.002, with
  $\rho = 0.025$; $\beta\in\{0, 0.5\}$, three values of $\tau$. The same
  $\eta$ values are used for both $\beta$. \\
Single-parameter $\rho$ & 1\,666 + 1\,500 &
  500 values of $\rho\in[0, 0.05]$~yr$^{-1}$, one in each stratum of width
  $10^{-4}$, with $\eta = 1.5$; $\beta\in\{0, 0.5\}$, three values of $\tau$.
  For $\beta = 0$, 166 of the $\tau = 0$ cells were run twice. \\
Low-$\rho$ extension & 340 &
  $\tau = 0$, $\beta = 0$; 340 equally spaced values of $\rho$ from $-0.029$
  to $0.004$~yr$^{-1}$, extending the single-parameter $\rho$ family to low
  consumption discount rates. \\
DICE comparison & 31 &
  $G(t)\equiv0$, $\tau = 0$, $\beta = 0$; 20 equally spaced values of $\rho$
  from $-0.015579$ to $0.024421$~yr$^{-1}$, $\rho\in\{0.03, 0.03616\}$, and
  nine values of $\rho$ chosen to give $r_\mu(2030)\approx$ 0.8, 1.0 (twice),
  1.2, 2, 3, 4, 4.7 and 5\%~yr$^{-1}$. \\
Infill & 2--484 per family &
  Additional $(\eta,\rho)$ samples for Fig.~5, targeted at bands of $r_\mu$,
  with $\eta\in[1,2]$ and $\rho$ down to $-0.041$~yr$^{-1}$; $\tau\in\{0.5,
  0.998\}$ or $\tau = 0.5$ only. Some were re-run from warm starts to move
  runs out of a local optimum; two runs used $s\in[0.2, 0.4]$. \\
\end{longtable}

% AUTHOR NOTE: the generator of the Latin hypercube samples is not archived.
% The design is inferred from the samples themselves: all 500 eta strata of
% width 0.002 are occupied, and 497 of 500 rho strata of width 1e-4 (the
% 3 empty strata are adjacent to values lying exactly on a stratum boundary,
% consistent with rounding to 6 decimals), whereas iid sampling would occupy
% about 316. Sample correlations of eta and rho: -5e-5 (beta = 0) and 0.05
% (beta = 0.5).
% AUTHOR NOTE: the warm-started re-runs used starting controls supplied by a
% batch driver that is not archived; their configuration files carry the
% standard initial guesses, so re-running them performs a fresh optimization.
% How the nine rho values targeting r_mu(2030) were found is not recorded.

%======================================================================
\section{Estimation of exogenous inputs}
\label{sec:estimation}
%======================================================================

The estimation procedures below were run once, offline; their results were
copied into the model configuration (Table~\ref{tab:provenance}).

\begin{table}[ht]
\centering
\caption{Provenance of the fitted inputs.}
\label{tab:provenance}
\small
\begin{tabular}{p{0.36\textwidth} l p{0.36\textwidth}}
\toprule
Model input & Source & Used as fitted? \\
\midrule
Lorenz exponents $p_k$, weights $\gamma_k$ & \ref{sec:pip}, \ref{sec:lorenz-fit} &
  weights in full; exponents rounded to 6 decimals \\
Gini trajectory $G(t)$ & \ref{sec:gini} & rounded to 5 decimals \\
$L(t)$, $\sigma(t)$ & \ref{sec:dice-fits} & rounded to 4--6 significant figures \\
$A(t)$ & \ref{sec:dice-fits} (closed form) & --- \\
$\theta_1$, $\xi$, $\Eland$, scalar parameters & \ref{sec:other-inputs} & --- \\
\bottomrule
\end{tabular}
\end{table}

\subsection{Global income distribution}
\label{sec:pip}

The global income distribution was constructed from the World Bank Poverty
and Inequality Platform~\cite{PIP} (downloaded 28 December 2025: 2\,556
survey rows for 172 countries, 1963--2024). Rows with complete decile shares,
population and per-capita GDP were retained, and for each country the most
recent survey year (1992--2024) was used. Each country--decile was treated as
a group containing one tenth of the country's population with income per
person
\begin{equation}
  y_{cd} = 10\,s_{cd}\,\mathrm{GDP}_c ,
\end{equation}
where $s_{cd}$ is the decile's share and $\mathrm{GDP}_c$ the per-capita GDP
reported by the platform. The 1\,720 groups were sorted by income; cumulative
population and income fractions, with the origin prepended, define the
empirical Lorenz curve. Its Gini coefficient, $G = 1 - 2\int_0^1\Lor\,\dd p$
by the trapezoidal rule, is 0.680860, for a population of
$7.719\times10^{9}$ and mean income \$11\,627.
% AUTHOR NOTE: (1) The platform mixes income and consumption surveys, and
% national, urban and rural coverage; nothing was filtered on either.
% (2) Units of GDP: the script describes reporting_gdp as "PPP $", but its
% magnitudes (US 2023: 65,505; India 2022: 2,065) look like constant US$ at
% market exchange rates. Check the PIP documentation before stating units.
% (3) For four countries with several rows in their latest year, the row used
% depended on an unstable sort: China 2022 national consumption, Haiti 2012
% income, Philippines 2023 income, St. Lucia 2015 consumption. Choosing
% consumption for Haiti and the Philippines would give G = 0.680706.

\subsection{Lorenz curve fit}
\label{sec:lorenz-fit}

The base Lorenz curve, Eq.~\eqref{eq:lorenz-base}, was fitted to the 1\,721
points $(x_j, y_j)$ of the empirical curve (including the origin) by
unweighted least squares, under the constraints $\gamma_k\ge0$ and
$\sum_k\gamma_k = 1$:
\begin{equation}
  \min_{p}\ \min_{\gamma}\ \sum_j\left(y_j - \sum_{k=0}^{3}\gamma_k x_j^{p_k}\right)^2 .
\end{equation}
For fixed exponents, the weights were found by sequential quadratic
programming (SLSQP~\cite{Kraft1988}; starting point $\gamma_k = 0.25$, function
tolerance $10^{-15}$, at most 2\,000 iterations, finite-difference gradients
with step $10^{-10}$). The exponents were optimized over the inner optimum
with L-BFGS-B~\cite{Byrd1995}, starting from $p = (1.1, 7.4, 13.7, 20)$
(equally spaced), with bounds $1.01\le p_k\le1000$ and tolerances $10^{-12}$
on the function and the projected gradient. The fit converged in 49
iterations to a residual sum of squares of $1.811\times10^{-3}$
($R^2 = 0.99997863$, maximum absolute error 0.0048), giving
$p = (1.50004,\ 4.36744,\ 14.0720,\ 135.060)$ and
$\gamma = (0.159871,\ 0.377619,\ 0.367125,\ 0.0953854)$. In the model the
exponents are rounded to six decimals and $\gamma_0$ is set to
$1 - \gamma_1 - \gamma_2 - \gamma_3$. The optimum is flat: other starting
points change the exponents by up to $1.1\times10^{-4}$ (relative) at the same
residual.

\subsection{Global Gini trajectory}
\label{sec:gini}

The trajectory $G(t)$ was derived from national projections of income decile
shares consistent with the Shared Socioeconomic Pathways~\cite{Narayan2023},
together with SSP projections of population and GDP per capita.
% AUTHOR NOTE: the decile-share file (ISO_level_projections_PC_model_projections.csv)
% is cited in the repository history as doi:10.1088/1748-9326/acbdb0; confirm
% that this is Narayan et al. (2023), and add the source of the SSP population
% and GDP files (pop_ssp_database.csv, All_GDP_percapita.csv), which is not
% documented.
For each SSP and for each of the 18 five-yearly years from 2015 to 2100, the
165 countries with complete data were split into deciles with population
$P_c/10$ and income per person $10\,s_{cd}(t)\,\mathrm{GDP}_c(t)$, sorted by
income, and the global Gini coefficient was computed with the trapezoidal rule
as above (here without prepending the origin). For SSP2 this gives
$G = 0.5962$ (2015), $0.5651$ (2020), $0.5244$ (2050) and $0.4647$ (2100).

The model's $G(t)$ is an unweighted least-squares fit of
Eq.~\eqref{eq:tf-exp} to all 18 SSP2 values, with $t_r = t - 2020$ (so 2015
enters as $t_r = -5$), which gives $a = 0.19517$, $b = -0.01058$~yr$^{-1}$
and $c = 0.37796$.
% AUTHOR NOTE: the exponential fit was not scripted; it was reconstructed and
% confirmed (scipy curve_fit, p0 = [0.2, -0.01, 0.4]). The fit's G(2020) =
% 0.57313 is 0.008 above the tabulated 0.565096.
% AUTHOR NOTE (data error): in both SSP input files, the rows for Taiwan
% (TWN) are identical to those for China (CHN), so China's population and
% income are counted twice, the second time with Taiwan's decile shares.
% Without TWN, the SSP2 Gini is 0.6063 (2020) and 0.4844 (2100), and the fit
% becomes a = 0.29188, b = -0.0075665, c = 0.31927. All published runs use the
% values with the duplicate. Decide how to report this.

\subsection{Productivity, population and carbon intensity}
\label{sec:dice-fits}

$A(t)$, $L(t)$ and $\sigma(t)$ follow DICE-2023~\cite{BarrageNordhaus2024},
whose recurrences run on five-year steps $k = 1, 2, \dots$, with
$t_r = 5(k-1)$. Each recurrence corresponds to a Gompertz function
(Eq.~\ref{eq:tf-gompertz}), exactly or approximately.

\paragraph{Population.} DICE's $L_{k+1} = L_k\,(L_\infty/L_k)^{0.145}$ is
exactly Gompertz, with $c = \ln(1-0.145)/5 = -0.031331$~yr$^{-1}$,
$L_0 = 7752.9$ million and $L_\infty = 10\,825$ million.

\paragraph{Carbon intensity.} DICE's $\sigma_{k+1} = \sigma_k
\exp(5g_k)$, with $g_k = -0.015\cdot0.96^{k-1}$, is exactly Gompertz with
$c = \ln(0.96)/5 = -0.0081644$~yr$^{-1}$ and
$\sigma_\infty = \sigma_0\,e^{5(-0.015)/(1-0.96)} = \sigma_0\,e^{-1.875}$, until
2155, when DICE's floor $g_k\ge-0.005$ starts to bind. The model uses the
Gompertz form throughout; it departs from DICE by about 5\% in 2210 and 25\% in
2280.

\paragraph{Productivity.} DICE's $A_{k+1} = A_k/(1-g_{A,k})$ with
$g_{A,k} = 0.066\,e^{-5\cdot0.0015(k-1)}$ is Gompertz in the limit of small
growth rates. The model uses that continuous approximation,
$c = -0.0015$~yr$^{-1}$ and
\begin{equation}
  A_\infty = A_0\exp\!\left(\frac{0.066}{5\times0.0015}\right) = A_0\,e^{8.8},
\end{equation}
with $A_0$ calibrated to DICE's 2020 gross output,
$A_0 = Q_0/(K_0^{\alpha}L_0^{1-\alpha}) = 739.619$ for $Q_0 = 135.7$ trillion
dollars, $K_0 = 295$ trillion dollars and $L_0 = 7.7529\times10^{9}$. Relative
to the DICE series, the model's $A(t)$ is lower by 1.45\% in 2050, 3.57\% in
2100 and 11.4\% in 2420.
% AUTHOR NOTE: this derivation of A(t) was reconstructed numerically; no script
% contains it. The fitting script in parameter_estimation/ fits A(t) directly
% (final value 5383526.78, c = -0.00154062, within 0.04% of DICE), but that fit
% was not used, although parameter_estimation/README.md says it was.

Least-squares fits of the Gompertz form to the DICE series (initial value
fixed; $A$ and $\sigma$ fitted in logarithms, $\sigma$ only to 2155) recover
these values: $c = -0.0313308$ for $L$ and $c = -0.0081644$,
$\sigma_\infty = 4.4680\times10^{-5}$ for $\sigma$.

\subsection{Other inputs}
\label{sec:other-inputs}

\begin{itemize}
  \item \emph{Land-use emissions}: $5.9\times10^9\,e^{-0.02107\,t_r}$
  tCO$_2$~yr$^{-1}$, DICE-2023's initial land-use emissions declining by 10\%
  per five-year period ($\ln(0.9)/5 = -0.02107$).
  \item \emph{Emission ratio}: $\xi = 1.4\,e^{-0.002085\,t_r}$, which falls to
  1.185 in 2100. (DICE-2023's ratio declines linearly from 1.40 in 2020 to
  1.21 in 2100 and is constant thereafter.)
  \item \emph{Backstop price}: a two-term exponential fit to the DICE-2023
  backstop price.
  \item \emph{Scalar parameters}: $\alpha$, $\delta$, $\psi_2$, $\theta_2$ and
  $K_0$ are DICE-2023 values.
  \item \emph{Temperature response}: the median model of
  Ricke and Caldeira~\cite{RickeCaldeira2014}.
\end{itemize}
% AUTHOR NOTE: (1) the source of the emission-ratio rate -0.002085 was not
% found; (2) no script or data for the backstop-price fit is archived;
% (3) no derivation of y_ref = 17210.79 was found (it affects results only
% through the clipping bounds, Sec. S5.9); (4) rho = 0.025 is labelled
% "DICE-2023 risk-adjusted rate" in the configurations, but DICE-2023's rate
% is exp(0.001 + 0.5*0.05) - 1 = 0.02634.

%======================================================================
\section{Construction of the figures}
\label{sec:figures}
%======================================================================

Throughout, the optimal carbon price in 2030 is the marginal abatement cost
(Section~\ref{sec:derived}) at $t = 2030$, and the warming in 2100 is
$\Delta T$ at $t = 2100$. Runs are matched to target values of $\tau$ or
$\beta$ within $\pm0.01$.

\subsection{Fig.~1: distribution of damage and tax burdens}

Fig.~1 is analytic. Income at rank $F$ is $y(F) = \Lor'(F)$ with $G = 0.57$,
evaluated on 2\,000 equally spaced ranks in $[10^{-6}, 1-10^{-6}]$. With damage
proportional to $y\cdot y^{-\beta}$, the cumulative damage share is
\begin{equation}
  D_\beta(F) = \frac{\int_0^F y^{1-\beta}\,\dd F'}{\int_0^1 y^{1-\beta}\,\dd F'},
\end{equation}
and the cumulative tax share is $T_\tau(F) = \int_0^F [y - c(y)]\,\dd F' /
\int_0^1 [y - c(y)]\,\dd F'$, with $c$ the schedule Eq.~\eqref{eq:schedule}
for $\eta = 1.5$ and $\Delta U$ chosen (Brent's method, tolerance $10^{-12}$)
so that the average tax rate is 20\%. The integrals are left Riemann sums on
the rank grid. Curves are shown for $\beta\in\{0, 0.5, 1\}$ and
$\tau\in\{0, 0.5, 0.998\}$.

\subsection{Fig.~2: reference cases}

Fig.~2 shows the carbon price and warming from 2025 to 2100 for the six
reference cases with $\beta\in\{0, 0.5\}$, $\tau\in\{0, 0.5, 0.998\}$. For each
$\tau$, the band between $\beta = 0$ and $\beta = 0.5$ is shaded.

\subsection{Fig.~3: dependence on $\tau$ and $\beta$}

The 2030 carbon price and 2100 warming of the 1\,681 contour-grid runs form a
$41\times41$ grid in $(\tau,\beta)$. For plotting, the grid values were
interpolated to $400\times400$ points with a piecewise cubic
(Clough--Tocher) interpolant on a Delaunay triangulation. Filled contours are
drawn at intervals of \$10~tCO$_2^{-1}$ from \$50 to \$350~tCO$_2^{-1}$ (values
clipped to that range, lines every \$20) and of 0.1~$^\circ$C from 1.5 to
3.5~$^\circ$C. The strip $0.998<\tau\le1$, which was not simulated, is shaded
grey. All runs are used, without a convergence filter.
% AUTHOR NOTE: Fig. 3 uses every contour run, whereas Fig. 4 drops runs that
% did not converge and Fig. 5 includes 33 points flagged as not converged.

\subsection{Fig.~4 and Fig.~S1: sensitivity to $\eta$ and $\rho$}

Fig.~4 uses the Monte Carlo and single-parameter families. A run is counted as
converged if the optimizer reported success with termination on the variable
or function tolerance; other runs are discarded. Where a parameter
combination was run twice, the converged run is kept, and between two
converged runs the one that used more evaluations. Every
$(\beta,\tau)$ cell of every family retains 500 runs.

For each family, $\beta$ and $\tau$, the distribution of the 2030 carbon price
and of the 2100 warming is shown as a Gaussian kernel density estimate with
Silverman's bandwidth~\cite{Silverman1986}, evaluated at 500 points over an
axis range shared across rows. A dashed line marks the central case, the run
closest to $(\eta,\rho) = (1.5, 0.025)$ in the metric
$\sqrt{(\eta-1.5)^2 + (\log_{10}\rho - \log_{10}0.025)^2}$. For the Monte
Carlo family with $\beta = 0$, this is $\eta = 1.509987$,
$\rho = 0.024683$, with 2030 carbon prices of \$42.79, \$71.11 and
\$144.96~tCO$_2^{-1}$ for $\tau = 0$, 0.5 and 0.998. Fig.~S1 plots the same
runs as the 2030 carbon price against the 2100 warming.

\subsection{The consumption discount rate $r_\mu$}
\label{sec:rmu}

Figs.~5 and~S3 are plotted against the consumption discount rate in 2030,
\begin{equation}
  r_\mu = \rho + \eta\sum_j \varpi_j \ln\frac{c_j(2031)}{c_j(2030)},
  \qquad
  \varpi_j = \frac{w_j\,c_j(2030)^{-\eta}}{\sum_k w_k\,c_k(2030)^{-\eta}},
  \label{eq:rmu}
\end{equation}
with $c_j = (1-s)\,y_{\mathrm{n},j}$: $\rho$ plus $\eta$ times the
marginal-utility-weighted logarithmic growth rate of consumption. For $G = 0$
it reduces to $\rho + \eta\ln[\bar c(2031)/\bar c(2030)]$.
% AUTHOR NOTE: the code that computed r_mu is not archived. Eq. (rmu) is
% reconstructed from the stored values: it reproduces them to 5e-5 percentage
% points for the base case and 0.002-0.004 pp across the Fig. 5 groups, better
% than simple (rather than logarithmic) growth (0.023 pp) or the discrete
% Ramsey rule of Eq. (ramsey) (9e-4 pp). Confirm before publication.

\subsection{Fig.~5: carbon price against $r_\mu$}

Fig.~5 combines the Monte Carlo, single-parameter and infill runs with
$1\le r_\mu\le6$\%~yr$^{-1}$ (5\,254 runs), grouped by $\beta$ and $\tau$.
Where the same parameters were run in more than one family, one run was kept;
this is how warm-started re-runs replace runs that had stopped in a local
optimum. The 2030 carbon price is shown on a logarithmic axis.
% AUTHOR NOTE: the curated point set (outputs/audits/points_cloud_1_6_balanced.csv)
% was built by a script that is not archived. Its selection rule and the
% "d_score" column are undocumented; 670 rows are marked as having duplicates
% in another family; 54 Monte Carlo runs with tau = 0.5 and r_mu of 5.1-5.5%
% that Fig. 4 counts as converged are absent ("balanced" thinning); and the
% 33 rows flagged as not converged are plotted.

\subsection{Fig.~S2: consumption growth across the distribution}

For the base case ($\tau = 0$, $\beta = 0$), the consumption growth rate of
node $j$ is $g_j(t) = 100\,[c_j(t+1)/c_j(t) - 1]$~\%~yr$^{-1}$, with
$c_j = (1-s)\,y_{\mathrm{n},j}$. In 2030 the population-mean growth rate is
$\bar g = \sum_j w_j g_j = 1.794$\%~yr$^{-1}$ and the
marginal-utility-weighted rate is $g_\mu = \sum_j\varpi_j g_j =
2.142$\%~yr$^{-1}$, with $\varpi_j$ as in Eq.~\eqref{eq:rmu}; the poorest node
grows at 2.779\% and the richest at 1.065\%~yr$^{-1}$. The contour panel
interpolates $g_j(t)$ for 2030--2100 (71 years $\times$ 32 nodes) to
$400\times400$ points by cubic interpolation, smooths it with a Gaussian
filter (standard deviations of about 2.1~yr and 2\% of rank), and draws
contours from 1.0 to 3.0\%~yr$^{-1}$.

\subsection{Fig.~S3: comparison with DICE-2023}

Fig.~S3 shows the 2030 carbon price against $r_\mu$ for (i) the 31 runs with
$G\equiv0$ (DICE comparison family), interpolated with PCHIP through the 30
runs with $r_\mu\le5.1$\%~yr$^{-1}$; (ii) the 840 runs with the empirical Gini
trajectory, $\tau = 0$ and $\beta = 0$ (single-parameter $\rho$ and low-$\rho$
families; where a run was duplicated, the one with the higher objective was
kept), interpolated with PCHIP through the 567 runs with
$0.5\le r_\mu\le5.5$\%~yr$^{-1}$ and evaluated at $r_\mu = 1, \dots, 5$\%;
and (iii) the DICE-2023 carbon prices in 2030 at constant discount rates of
1--5\% (\$595, 224, 114, 66 and 43~tCO$_2^{-1}$, from the DICE-2023
workbook~\cite{BarrageNordhaus2024}).
Table~\ref{tab:figs3} gives the plotted values.

\begin{table}[ht]
\centering
\caption{Optimal 2030 carbon price (\$ tCO$_2^{-1}$) against the consumption
discount rate $r_\mu$, as plotted in Fig.~S3.}
\label{tab:figs3}
\small
\begin{tabular}{l r r r r r}
\toprule
$r_\mu$ (\%~yr$^{-1}$) & 1 & 2 & 3 & 4 & 5 \\
\midrule
This model, $G = 0$ & 396.40 & 213.99 & 112.12 & 62.33 & 38.43 \\
This model, empirical Gini & 650.94 & 380.95 & 204.74 & 107.17 & 59.65 \\
DICE-2023 & 595 & 224 & 114 & 66 & 43 \\
\bottomrule
\end{tabular}
\end{table}

%======================================================================
\section{Verification}
\label{sec:verification}
%======================================================================

The repository contains a regression suite that compares the model with a
baseline computed by the code used for the manuscript: forward replays of the
nine reference runs with their published controls (relative tolerance
$10^{-12}$ for $\tau = 0$ and $10^{-9}$ for $\tau>0$, set by the tax root
solve), the published results files themselves, a DICE-comparison run, six
cases exercising the options not used in the manuscript, three short
optimizations, and the individual numerical kernels.

Table~\ref{tab:checkpoints} gives reference values from the replay of the
base case ($\tau = 0$, $\beta = 0$; $J = -6.487938701598822\times10^{9}$).
For $\tau = 0.5$, $\beta = 0.5$, $J = -6.508855233019977\times10^{9}$; for
$\tau = 0.998$, $\beta = 1$, $J = -6.510344441554529\times10^{9}$.

\begin{table}[ht]
\centering
\caption{Reference values from the replay of the base case
($\tau = 0$, $\beta = 0$).}
\label{tab:checkpoints}
\small
\begin{tabular}{l l l l}
\toprule
Quantity & 2020 & 2030 & 2100 \\
\midrule
$f$ & 1.543376572092825 & 1.629652538702295 & 2.395393167065027 \\
$s$ & 0.4786271018324223 & 0.2296397516562726 & 0.2300201664610637 \\
$K$ (\$) & $2.95\times10^{14}$ & $3.402759806869976\times10^{14}$ &
  $1.343287212994556\times10^{15}$ \\
$\Delta T$ ($^\circ$C) & 1.16966 & 1.470360532030760 & 3.107513349030306 \\
$\yn$ (\$) & $1.741645953981179\times10^{4}$ & $2.010579385432332\times10^{4}$ &
  $6.519107748253394\times10^{4}$ \\
$\mu$ & $3.996318342682483\times10^{-2}$ & $1.280010855705555\times10^{-1}$ &
  $6.313729676087462\times10^{-1}$ \\
$\Omega$ & $4.743218355585199\times10^{-3}$ & $7.495515646431158\times10^{-3}$ &
  $3.347956815633018\times10^{-2}$ \\
$U$ & $-2.360533615801890\times10^{8}$ & $-1.915743595569843\times10^{8}$ &
  $-1.180399699134512\times10^{8}$ \\
\bottomrule
\end{tabular}
\end{table}

%======================================================================
\section{Numerical properties of the method}
\label{sec:caveats}
%======================================================================

The following properties are part of the method as used for the published
results.
\begin{enumerate}
  \item \emph{Lagged damage distribution.} Each step applies the previous
  step's distribution of damage across incomes, scaled to the current
  aggregate damage (Section~\ref{sec:damage}). The first step uses uniform
  damage.
  \item \emph{Mean pre-tax income.} The node incomes $y_{\mathrm{pre},i}$
  average to $\yd$ only approximately (Section~\ref{sec:step}); the closed form
  of the flat tax assumes that they do.
  \item \emph{Quadrature error.} With 32 nodes, the mean of the empirical
  Lorenz density is in error by $3$--$7\times10^{-7}$ (relative).
  \item \emph{Safeguards.} Damage fractions are confined to
  $[0, 1-\varepsilon]$, abatement expenditure to at most $(1-\varepsilon)\yd$,
  consumption to at least $\varepsilon$ in the utility, and the flat tax rate
  to at most $1-\varepsilon$, with $\varepsilon = 10^{-12}$ in absolute terms
  for quantities of different scale (fractions and dollars per person). The
  flat-tax case is selected by $|\etat - 1|<\varepsilon$, whereas logarithmic
  utility is used only for $\eta = 1$ exactly.
  \item \emph{Infeasible trajectories} receive the values $-10^{17}$ or
  $-10^{30}$ at the optimizer interface (Section~\ref{sec:nlopt}).
  \item \emph{State constraints.} $\Ecum\ge0$ is the only constraint on the
  state; $\mu$ is effectively unbounded ($\mu_{\max} = 10^{20}$).
\end{enumerate}

%======================================================================
\begin{thebibliography}{99}
%======================================================================

\bibitem{BarrageNordhaus2024}
L.~Barrage, W.~Nordhaus, Policies, projections, and the social cost of carbon:
Results from the DICE-2023 model. \emph{Proc. Natl. Acad. Sci. U.S.A.}
\textbf{121}, e2312030121 (2024).

\bibitem{Caldeira2023}
K.~Caldeira \emph{et al.}, \emph{Environ. Res. Lett.} (2023),
doi:10.1088/1748-9326/acf949.
% AUTHOR NOTE: complete this reference (authors, title, volume).

\bibitem{RickeCaldeira2014}
K.~L. Ricke, K.~Caldeira, Maximum warming occurs about one decade after a
carbon dioxide emission. \emph{Environ. Res. Lett.} \textbf{9}, 124002 (2014).

\bibitem{Harris2020}
C.~R. Harris \emph{et al.}, Array programming with NumPy. \emph{Nature}
\textbf{585}, 357--362 (2020).

\bibitem{Virtanen2020}
P.~Virtanen \emph{et al.}, SciPy 1.0: fundamental algorithms for scientific
computing in Python. \emph{Nat. Methods} \textbf{17}, 261--272 (2020).

\bibitem{NLopt}
S.~G. Johnson, The NLopt nonlinear-optimization package,
\url{https://github.com/stevengj/nlopt}.

\bibitem{Rowan1990}
T.~H. Rowan, Functional stability analysis of numerical algorithms. PhD thesis,
University of Texas at Austin (1990).

\bibitem{Brent1973}
R.~P. Brent, \emph{Algorithms for Minimization without Derivatives}
(Prentice-Hall, 1973).

\bibitem{FritschCarlson1980}
F.~N. Fritsch, R.~E. Carlson, Monotone piecewise cubic interpolation.
\emph{SIAM J. Numer. Anal.} \textbf{17}, 238--246 (1980).

\bibitem{FritschButland1984}
F.~N. Fritsch, J.~Butland, A method for constructing local monotone piecewise
cubic interpolants. \emph{SIAM J. Sci. Stat. Comput.} \textbf{5}, 300--304
(1984).

\bibitem{McKay1979}
M.~D. McKay, R.~J. Beckman, W.~J. Conover, A comparison of three methods for
selecting values of input variables in the analysis of output from a computer
code. \emph{Technometrics} \textbf{21}, 239--245 (1979).

\bibitem{PIP}
World Bank, Poverty and Inequality Platform, \url{https://pip.worldbank.org}
(accessed 28 December 2025).

\bibitem{Kraft1988}
D.~Kraft, A software package for sequential quadratic programming. Tech. Rep.
DFVLR-FB 88-28, Deutsche Forschungs- und Versuchsanstalt f\"ur Luft- und Raumfahrt (1988).

\bibitem{Byrd1995}
R.~H. Byrd, P.~Lu, J.~Nocedal, C.~Zhu, A limited memory algorithm for bound
constrained optimization. \emph{SIAM J. Sci. Comput.} \textbf{16}, 1190--1208
(1995).

\bibitem{Narayan2023}
K.~B. Narayan, B.~C. O'Neill, S.~T. Waldhoff, C.~Tebaldi, Non-parametric
projections of national income distribution consistent with the Shared
Socioeconomic Pathways. \emph{Environ. Res. Lett.} \textbf{18}, 044013 (2023).
% AUTHOR NOTE: verify this reference against doi:10.1088/1748-9326/acbdb0.

\bibitem{Silverman1986}
B.~W. Silverman, \emph{Density Estimation for Statistics and Data Analysis}
(Chapman and Hall, 1986).

\end{thebibliography}

\end{document}
